% Generated draft: requires mathematical and source review.
\documentclass[11pt]{article}
\usepackage[margin=25mm]{geometry}
\usepackage{fontspec}
\setmainfont{Latin Modern Roman}
\usepackage{amsmath,amssymb,mathtools,unicode-math}
\setmathfont{Latin Modern Math}
\usepackage{enumitem,needspace,xurl,hyperref}
\setlength{\parindent}{0pt}
\setlength{\parskip}{5pt}
\setlength{\emergencystretch}{3em}
\newcommand{\N}{\mathbb{N}}
\newcommand{\Z}{\mathbb{Z}}
\newcommand{\Q}{\mathbb{Q}}
\newcommand{\R}{\mathbb{R}}
\newcommand{\C}{\mathbb{C}}
\newcommand{\F}{\mathbb{F}}
\newcommand{\A}{\mathbb{A}}
\newcommand{\PP}{\mathbb{P}}
\newcommand{\E}{\mathbb{E}}
\newcommand{\eps}{\varepsilon}
\newcommand{\dd}{\,\mathrm{d}}
\begin{document}
\section*{UnsolvedMath ID: 30001303}
\textbf{Source problem code:} catalogue entry
\subsection{Definitions and mathematical statement}
Let \(K^{\mathrm{ab}}\) denote the maximal abelian extension of the field \(K\) in a fixed separable closure. The supplied record does not specify \(K\), the admissible values of \(r\), the variety \(X_r\), the object \(A\), the map \(\Phi_C\), or the construction and hypotheses defining a generalized Heegner cycle \(\Delta_\varphi\). It also does not define the subscript \(0\) in \(CH^{r+1}(X_r)_0(H)\). These symbols therefore retain their source-dependent meanings; the following is a conditional transcription requiring those definitions before it becomes a fully specified conjecture.

The notation \(CH^{r+1}\) denotes a Chow group of codimension-\(r+1\) cycles modulo rational equivalence. The record does not specify its coefficient convention or precisely how the notation \({}_0(H)\) is formed. In the asserted target, \(A(H)\) denotes the group of \(H\)-rational points of the intended object \(A\), and \(A(H)\otimes_{\Z}\Q\) denotes its rationalization, provided the required commutative group structure has been specified.

\textbf{Conjecture as supplied.} For every field \(H\subseteq K^{\mathrm{ab}}\) and every generalized Heegner cycle
\[
\Delta_\varphi\in CH^{r+1}(X_r)_0(H),
\]
the source asks whether
\[
\Phi_C(\Delta_\varphi)\in A(H)\otimes_{\Z}\Q
\]
and
\[
\Phi_C(\Delta_\varphi^\sigma)
 =\Phi_C(\Delta_\varphi)^\sigma
\qquad\text{for every }\sigma\in\operatorname{Gal}(H/K).
\]
For the displayed Galois group to have its stated meaning, \(H\) must contain \(K\); thus the intended field condition appears to be \(K\subseteq H\subseteq K^{\mathrm{ab}}\). This containment is implicit rather than explicit in the supplied statement. No finiteness condition on \(H/K\) is supplied. The constructions of both Galois actions and the identification needed to interpret the displayed membership in the target of \(\Phi_C\) remain unspecified.
\subsection{Short English statement}
The conjecture asks whether the image of every generalized Heegner cycle defined over an abelian extension \(H\) of \(K\) is an \(H\)-rational point of the intended group \(A\), with rational coefficients, and whether taking this image commutes with every element of \(\operatorname{Gal}(H/K)\). The source omits definitions and hypotheses needed to make this a fully specified mathematical statement.
\subsection{Sources}
\begin{itemize}
\item Supplied local UnsolvedMath registry record, permanent ID \(30001303\), title: Galois Equivariance of Generalized Heegner-Cycle Images; source problem code: \texttt{catalogue entry}. Collection: Oberwolfach Reports Open Problems. The registry records proposed year 2009 and supplies no named proposer.
\item Supplied source citation: \emph{Explicit Methods in Number Theory (2010)}, Oberwolfach Reports, pp.\ 1843--1920, DOI \texttt{10.4171/owr/2009/33}. \url{https://doi.org/10.4171/owr/2009/33}. The background describes these reports as documenting mathematical workshops and their open-problem discussions.
\item Supplied literature source, labeled ``Darmon Oberwolfach report'': \url{https://ems.press/content/serial-article-files/46233?nt=1}. This label alone does not establish the proposer of this conjecture.
\item Supplied literature source, labeled ``Generalized Heegner cycles theorem'': \url{https://doi.org/10.1215/00127094-2142056}. The record does not provide a theorem number or its hypotheses.
\item Literature assessment dated exactly \texttt{2026-08-21}: the supplied curation labels the problem partially solved and states that the complex algebraicity conjecture is not known in full generality, while \(p\)-adic analogues establish predicted algebraicity properties in broad settings. This is supplied dated curation derived from the OWR material, not an independently verified literature assessment. Its time-sensitive status and cited claims require independent verification.
\item The registry supplies the statement label \texttt{corrected\_verified} and the verification description that the field-of-definition and Galois-equivariance conjecture is complete and its notation normalized. These are supplied registry assertions, not independent verification of this draft.
\end{itemize}
\end{document}
